\cotefile{157-2}
\scanpage{131}
\mspage{76}
% archive sheet number: 130

\diagram{\begin{tikzcd}[ampersand replacement=\&,column sep=0.45em,row sep=0.68em,font=\scriptsize]
\&\&\&\underline{\mathrm{Fibhot}}\,S
 \arrow[dd,hook,"\unsure{\delta|d}"]
 \arrow[ddddlll,dashed,bend right=12,"\Phi_0"']
 \arrow[ddddrrr,dashed,bend left=12,"\Psi_0"]\\
\&\&\underline{\mathrm{Liss}}_0\,S \arrow[ur,hook] \arrow[ddd,hook]\&\&
 \underline{\mathrm{Prop}}_0\,S \arrow[ul,hook] \arrow[ddd,hook]\\
\&\underline{\mathrm{Fib}}_0\,S \arrow[ur,hook] \arrow[ddd,hook]\&
 \&\mathrm{Cat}_{/S} \arrow[d,"C"] \arrow[ddddlll,dashed,bend right=9,"\Phi"'] \arrow[ddddrrr,dashed,bend left=9,"\Psi"]\&\&
 \underline{\mathrm{Cofib}}_0\,S \arrow[ul,hook] \arrow[ddd,hook]\\
\&\&\&\underline{\mathrm{Parf}}\,S \arrow[uul,hook] \arrow[uur,hook]
 \arrow[ull,dotted,leftarrow,"\deleted{\unsure{\triangleleft}}"'{pos=0.94}] \arrow[urr,dotted,leftarrow,"\deleted{\unsure{\triangleright}}"{pos=0.94}]
 \arrow[ddd,dashed,"\mathrm{id}"]\\
\underline{\mathrm{Fibsc}}_0\,S \arrow[uur] \arrow[dd,hook]\&\&
 \underline{\mathrm{Liss}}\,S \arrow[uur,hook]\&\&
 \underline{\mathrm{Prop}}\,S \arrow[uul,hook]\&\&
 \underline{\mathrm{Cofibsc}}_0\,S \arrow[uul] \arrow[dd,hook]\\
\&\underline{\mathrm{Fib}}\,S \arrow[ur,hook]\&\&\&\&
 \underline{\mathrm{Cofib}}\,S \arrow[ul,hook]\\
\underline{\mathrm{Fibsc}}\,S \arrow[ur]\&\&\&
 \underline{\mathrm{Parf}}\,S \arrow[uul,dashed,hook] \arrow[uur,dashed,hook]\&\&\&
 \underline{\mathrm{Cofibsc}}\,S \arrow[ul] \arrow[ul,phantom,"\deleted{\unsure{\downarrow}}"']
\end{tikzcd}}

J'ai \unsure{tracé|dessiné} en \inserted{pointillés} verts les cinq foncteurs $\Phi,\Psi,\Phi_0,\Psi_0,C$ qui \deleted{\unsure{par passage au dérivé}} \inserted{par passage aux localisés} donnent des quasi-inverses pour toutes les équivalences de catégories énoncées dans \textcircled{I} \textcircled{II} du présent \unsure{Scholie}. \deleted{\unsure{De}} Pour préciser, le foncteur

\[
\Phi:\mathrm{Cat}_{/S}\longrightarrow\underline{\mathrm{Fibsc}}\,S
\]

associe à tout $X$ de $\mathrm{Cat}_{/S}$ la catégorie scindée sur $S$

\[
s\longmapsto {}_s\backslash X
\]

sur $S$. Ce foncteur, et ses restrictions à $\underline{\mathrm{Liss}}\,S$, $\underline{\mathrm{Fib}}\,S$ et, plus généralement, à toute sous-catégorie pleine $\unsure{\mathcal{F}|F}$ de $\mathrm{Cat}_{/S}$ contenant $\underline{\mathrm{Fib}}\,S$, \deleted{établit} est \unsure{compatible} avec ces localiseurs $W_S$ d'une part (sur $\mathcal{F}$ et $\underline{\mathrm{Fibsc}}\,S$), $W_S^d$ d'autre part, et passe \inserted{donc} aux localisés correspondants, pour fournir le quasi-inverse \inserted{cherché} de

\[
(\underline{\mathrm{Fibsc}}\,S)\Sigma^{-1}\longrightarrow\mathcal{F}\,\Sigma^{-1}
\qquad(\Sigma=W_S\text{ ou }W_S^d).
\]

On peut aussi, pour tte sous-catégorie pleine $\mathcal{F}'$ intermédiaire

\[
\underline{\mathrm{Fib}}_S\subset\mathcal{F}'\subset\mathcal{F}\subset\mathrm{Cat}_{/S}
\]

regarder $\Phi$ comme un foncteur $\mathcal{F}\longrightarrow\mathcal{F}'$ (plus précisément, prendre le composé $\mathcal{F}\xrightarrow{\Phi|\mathcal{F}}\underline{\mathrm{Fibsc}}\,S\longrightarrow\underline{\mathrm{Fib}}\,S\hookrightarrow\mathcal{F}'$) soit $\Phi_{\mathcal{F}',\mathcal{F}}$, et \deleted{\unsure{de ce qui}} ce qui vient d'être dit résulte que les deux foncteurs

\diagram{\begin{tikzcd}[ampersand replacement=\&,column sep=large]
\mathcal{F}'\arrow[r,hook]\&\mathcal{F}\arrow[l,bend left=16,"\Phi_{\mathcal{F}',\mathcal{F}}"]
\end{tikzcd}}
\scanpage{132}
\mspage{77}
% archive sheet number: 131

établissent des équivalences quasi-inverses l'une de l'autre entre

\diagram{\begin{tikzcd}[ampersand replacement=\&,column sep=large]
{\unsure{\mathcal{F}|F}'\Sigma^{-1}}\arrow[r,"\approx"]\&\mathcal{F}\Sigma^{-1}\arrow[l,bend left=15,"{\underset{\approx}{\Phi_{\mathcal{F}',\mathcal{F}}}}"]\&(\Sigma=W_S\text{ ou }W_S^d).
\end{tikzcd}}

Énoncé dual pour le foncteur

\[
\Psi:\mathrm{Cat}_{/S}\longrightarrow\underline{\mathrm{Cofibsc}}\,S,
\]

associant à $X$ de $\mathrm{Cat}_{/S}$ la catégorie cofibrée scindée sur $S$

\[
s\longmapsto X_{/s}.
\]

Il donne des \unsure{couples} de foncteurs quasi-inverses

\diagram{\begin{tikzcd}[ampersand replacement=\&,column sep=large]
(\underline{\mathrm{Cofibsc}}\,S)\Sigma^{-1}\arrow[r]\&\mathcal{F}\Sigma^{-1}\arrow[l,bend left=14,"\Psi_{\mathcal{F}}"]\&(\Sigma=W_S\text{ ou }W_S^g),
\end{tikzcd}}

pour toute sous-catégorie pleine \inserted{$\mathcal{F}$} de $\mathrm{Cat}_{/S}$ telle que

\[
\underline{\mathrm{Cofib}}_S\subset\mathcal{F}\subset\mathrm{Cat}_{/S},
\]

et des couples de foncteurs quasi-inverses

\diagram{\begin{tikzcd}[ampersand replacement=\&,column sep=large]
\mathcal{F}'\Sigma^{-1}\arrow[r,"\approx"]\&\mathcal{F}\Sigma^{-1}\arrow[l,bend left=15,"\approx"',"{\Psi_{\mathcal{F}',\mathcal{F}}}"]\&(\Sigma=W_S\text{ ou }W_S^g)
\end{tikzcd}}

pour deux sous-catégories pleines $\mathcal{F}',\mathcal{F}$

\[
\underline{\mathrm{Cofib}}_S\subset\mathcal{F}'\subset\mathcal{F}\subset\mathrm{Cat}_{/S}.
\]

Les foncteurs $\Phi,\Psi$ \deleted{induisent sur} \inserted{envoient} la sous-catégorie pleine $\underline{\mathrm{Fibhot}}\,S$ \deleted{\unsure{des}} ``l'étage'', dans les sous-catégories $\underline{\mathrm{Fibsc}}_0\,S$, $\underline{\mathrm{Cofibsc}}_0\,S$ respectivement, donc induisent à l'étage des foncteurs $\Phi_0,\Psi_0$

\diagram{\begin{tikzcd}[ampersand replacement=\&,column sep=large]
\deleted{\unsure{\Phi_0}}\underline{\mathrm{Fibsc}}_0\,S\&\underline{\mathrm{Fibhot}}\,S\arrow[l,"\Phi_0"']\arrow[r,"\Psi_0"]\&\underline{\mathrm{Cofibsc}}_0\,S,
\end{tikzcd}}

jouant à l'étage exactement le $=$ rôle que $\Phi,\Psi$ \unsure{au|en} rez de chaussée. La seule différence, c'est qu'à l'étage on a $W_S=W_S^d=W_S^g\deleted{\unsure{=W_S}}\,(=W_S^{\unsure{f|t}})$,
\scanpage{133}
\mspage{78}
% archive sheet number: 132

donc il n'y a qu'une sorte de \inserted{\unsure{cat|cas}} \unsure{localisées|localisés} à examiner. On trouve ainsi

\diagram{\begin{tikzcd}[ampersand replacement=\&,column sep=large]
(\underline{\mathrm{Fibsc}}_0\,S)\deleted{\unsure{\Sigma}}W_S^{-1}\arrow[r,"\approx"]\&{\unsure{\mathcal{F}|F}W_S^{-1}}\arrow[l,bend left=16,"\approx"',"{\Phi_{0\mathcal F}}"]
\end{tikzcd}}

\diagram{\begin{tikzcd}[ampersand replacement=\&,column sep=large]
\mathcal{F}'\deleted{\unsure{\Sigma}}W_S^{-1}\arrow[r,"\approx"]\&\mathcal{F}\deleted{\unsure{\Sigma}}W_S^{-1}\arrow[l,bend left=16,"\approx"',"{\Phi_{0\mathcal{F}',\mathcal{F}}}"]
\end{tikzcd}}

pour $\mathcal{F}$ (ou $\mathcal{F},\mathcal{F}'$) sous-catégorie(s) pleine(s) de $\deleted{\unsure{\mathcal{F}}}\underline{\mathrm{Fibhot}}\,S$ telle que

\[
\underline{\mathrm{Fib}}_0\,S\subset\mathcal{F}\subset\underline{\mathrm{Fibhot}}\,S
\qquad\text{resp.}\qquad
\underline{\mathrm{Fib}}_0\,S\subset\mathcal{F}'\subset\mathcal{F}\subset\underline{\mathrm{Fibhot}}\,S,
\]

et de même du côté droit pour les sous-catégories \inserted{pl.} $\mathcal{F}$ ou $\mathcal{F},\mathcal{F}'$ de $\underline{\mathrm{Fibhot}}\,S$ satisfaisant

\[
\underline{\mathrm{Cofib}}_0\,S\subset\mathcal F\subset\underline{\mathrm{Fibhot}}\,S
\qquad\text{respt.}\qquad
\underline{\mathrm{Cofib}}_0\,S\subset\mathcal F'\subset\mathcal F\subset\underline{\mathrm{Fibhot}}\,S.
\]

\inserted{\uline{NB} La catégorie $\underline{\mathrm{Parf}}\,S$ contient à la fois $\underline{\mathrm{Fib}}_0\,S$ et $\underline{\mathrm{Cofib}}_0\,S$, et fait donc partie des catégories admissibles \unsure{dont varie|dont l'une est} $\mathcal{F}$, ou $\mathcal{F}'$.}

Ainsi, les quatre foncteurs $\Phi,\Psi,\Phi_0,\Psi_0$ donnent tous les quasi-inverses qu'on peut souhaiter au \deleted{\unsure{premier étage}} \inserted{rez de chaussée} d'une part ($\Phi$ et $\Psi$), à l'étage de l'autre ($\Phi_0,\Psi_0$) \deleted{\unsure{(ces catégories, cf. ci-dessus)}}. Ils ramènent \unsure{tous les localisés pertinents|toutes les localisées pertinentes} aux \deleted{\unsure{trois}} \inserted{quatre} suivants

\[
\begin{aligned}
&(\mathrm{Cat}_{/S})W_S^{-1},(\mathrm{Cat}_{/S})(W_S^d)^{-1},(\mathrm{Cat}_{/S})(W_S^g)^{-1}\\
&\hspace{11em}\text{(au rez de chaussée)}\\
&(\underline{\mathrm{Fibhot}}\,S)W_S^{-1}\qquad\text{(à l'étage),}
\end{aligned}
\]

\textcircled{VII} \deleted{\unsure{après un choix|ap. un choix} \unsure{principe|privilégié} \unsure{trouvés|trois} \unsure{pleins|trois} à l'étage,} Il reste à comparer le localisé à l'étage avec les localisés \deleted{\unsure{grossiers}} au rez de chaussée. Les flèches verticales \deleted{\unsure{gauches}} descendant de l'étage au rez de chaussée sont compatibles avec $W_S^d$ du côté gauche du diagramme, avec $W_S^g$ du côté droit, donc elles induisent \inserted{essentiellement} \uline{deux} foncteurs différents
\scanpage{134}
\mspage{79}
% archive sheet number: 133

\inserted{$i,j$} \unsure{sur|de} $(\underline{\mathrm{Fibhot}}\,S)W_S^{-1}$, dans $\mathrm{Cat}_{/S}(W_S^d)^{-1}$ et $(\mathrm{Cat}_{/S})(W_S^g)^{-1}$:

\diagram{\begin{tikzcd}[ampersand replacement=\&,column sep=0.35em,row sep=2.3em,font=\small]
\&(\underline{\mathrm{Fibhot}}\,S)W_S^{-1}\arrow[dl,"i"']\arrow[dr,"j"]\arrow[dd,dashed,"\delta","\approx"']\&\\
{\mathrm{HOT}_{W}(S)\overset{\mathrm{def}}{\unsure{\cong|\simeq|\approx}}\mathrm{Cat}_{/S}(W_S^d)^{-1}}\arrow[dr,dashed,"\rho"']\&\&
{\mathrm{Cat}_{/S}(W_S^g)^{-1}\overset{\mathrm{def}}{\approx}\mathrm{HOT}_{W}^{\unsure{\circ|o}}(S)}\arrow[dl,dashed,"\sigma"]\\
\&{\begin{gathered}(\mathrm{Cat}_{/S})W_S^{-1}\deleted{\unsure{\mathrm{hot}}}\\
\overset{\mathrm{def}}{\approx}\mathrm{HOT}_{W}^{\unsure{\mathrm{lc}|\mathrm{le}}}(S)\end{gathered}}\&,
\end{tikzcd}}

de façon à faire commuter le carré complété par les pointillés, représentant les foncteurs de localisation $\rho,\sigma$. On a bien sûr \unsure{q}

\[
\begin{aligned}
\rho i=\sigma j:\;(\underline{\mathrm{Fibhot}}\,S)W_S^{-1}&\longrightarrow(\mathrm{Cat}_{/S})W_S^{-1}\\
&\text{induit par l'inclusion}\\
&\delta:\underline{\mathrm{Fibhot}}\,S\hookrightarrow\mathrm{Cat}_{/S}.
\end{aligned}\tag*{(*)}
\]

Nous avons énoncé \deleted{que} dans \textcircled{I} que cette flèche est une équivalence, de sorte que $\rho$ et $\sigma$ peuvent être interprétés, à équiv. \inserted{de foncteurs} près, comme des rétractions de $\mathrm{HOT}_{W}(S)$ et de $\mathrm{H\unsure{O|\ell}T}_{W}^{\circ}(S)$ sur \unsure{leur|une} sous-catégorie commune, s'interprétant soit comme $(\underline{\mathrm{Fibhot}}\,S)W_S^{-1}$, soit comme $(\mathrm{Cat}_{/S})W_S^{-1}$.

Il reste à décrire le \unsure{\deleted{foncteur}|foncteur} quasi-inverse de ($*$), ce qui donnera en même temps les quasi-inverses \deleted{\unsure{à gauche (dont, p. ex., soit $f$, soit $g$)}} de tous les foncteurs verticaux (de l'étage vers le rez de chaussée) dans le diagramme \unsure{préc\textsuperscript{dt}|p.~76}, plus exactement pour leurs localisés pour $W_S^{\deleted{\unsure{d}}}$ \deleted{\unsure{resp.} $\unsure{\Rightarrow}$ \unsure{à des} $\unsure{\Rightarrow}$ \unsure{et ceci pour} $W_S^{\unsure{d}}$ \unsure{(de l'étage)}, et}
\scanpage{135}
\mspage{80}
% archive sheet number: 134

C'est ce qui est obtenu à l'aide du foncteur ``montant'' (du \unsure{r. de ch.} vers l'étage) $C$, qu'on peut regarder comme un foncteur $\mathrm{Cat}_{/S}\unsure{\mathrel{\relbar\joinrel\relbar\joinrel\vert}|\longrightarrow}\underline{\mathrm{Fibhot}}\,S$, mais qui en fait prend ses valeurs dans $\underline{\mathrm{Parf}}\,S$. (C'est là un détail sans grande importance, me semble-t-il.)

\[
\mathrm{Cat}_{/S}\xrightarrow{C}\underline{\mathrm{Parf}}\,S\subset\underline{\mathrm{Fibhot}}\,S.
\]

Ce foncteur \deleted{\unsure{transforme}} est compatible avec les localiseurs $W_S$ des deux côtés, \deleted{\unsure{donc}} \inserted{et} établit des équivalences (quasi-inverses des équivalences

\[
(\underline{\mathrm{Parf}}\,S)\deleted{\unsure{\Sigma}}W_S^{-1}\xrightarrow{\approx}(\mathrm{Cat}_{/S})W_S^{-1}
\quad\text{et}\quad
(\underline{\mathrm{Fibhot}}\,S)W_S^{-1}\overset{\approx}{\unsure{\mathrel{\relbar\joinrel\relbar\joinrel\vert}|\longrightarrow}}(\mathrm{Cat}_{/S})W_S^{-1}\bigr)
\]

\[
(\mathrm{Cat}_{/S})W_S^{-1}\xrightarrow{\approx}(\underline{\mathrm{Parf}}\,S)W_S^{-1}
\bigl(\overset{\approx}{\unsure{\mathrel{\relbar\joinrel\relbar\joinrel\vert}|\longrightarrow}}(\underline{\mathrm{Fibhot}}\,S)W_S^{-1}\bigr).
\]

Donc, quand les trois catégories $\mathrm{HOT}$ pour $S$ sont interprétées en termes de localisées de $\mathrm{Cat}_{/S}$, c'est le foncteur $C$ \inserted{(plus \unsure{exactemt|exactement} le composé $\mathrm{Cat}_{/S}\xrightarrow{C}\underline{\mathrm{Parf}}\hookrightarrow\mathrm{Cat}_{/S}$)} qui fournit à la fois le foncteur d'inclusion (!) $i$, et le foncteur d'inclusion $j$ dans $\mathrm{HOT}_{W}(S)$ et $\mathrm{HOT}_{W}^{\unsure{\circ|o}}(S)$ respectivement. Le foncteur en question

\[
C':\mathrm{Cat}_{/S}\longrightarrow\mathrm{Cat}_{/S}
\]

satisfait bien

\[
C'(W_S)\subset W_S^d\cap W_S^g
\]

et définit les deux foncteurs

\diagram{\begin{tikzcd}[ampersand replacement=\&,row sep=0.7em,column sep=large]
\&(\mathrm{Cat}_{/S})(W_S^d)^{-1}\approx\mathrm{HOT}_W(S)\\
{\begin{gathered}(\mathrm{Cat}_{/S})W_S^{-1}\\\approx\mathrm{HOT}_{W}^{\unsure{\mathrm{lc}|\mathrm{le}}}(S)\end{gathered}}\arrow[ur,"i"]\arrow[dr,dash,"j"']\&\\
\&(\mathrm{Cat}_{/\unsure{J|S}})(W_S^g)^{-1}\approx\mathrm{HOT}_{W}^{\circ}(S),
\end{tikzcd}}

Cela nous permet d'expliciter les formules d'adjonction de \textcircled{IV}, quand on y interprète $\unsure{\xi|x},\unsure{y|\eta}$
\scanpage{136}
\mspage{81}
% archive sheet number: 135

comme \inserted{\unsure{(,)}} provenant de $X,Y$ de $\mathrm{Cat}_{/S}$:

\[
\boxed{\mathrm{Hom}_{W_S^d}(\deleted{\unsure{C}}X,CY)\simeq\mathrm{Hom}_{W_S^g}(X,CY)\simeq\mathrm{Hom}_{W_S}(X,Y),}
\]

% slanted
\margin{Cette formule suggère que si $X\in\mathrm{Ob}\,\mathrm{Cat}_{/S}$, $Y\in\mathrm{Ob}\,\underline{\mathrm{Fibhot}}\,S$, alors $\mathrm{Hom}_{W_S^d}(X,Y)\unsure{\simeq_{?}|\simeq}\mathrm{Hom}_{W_S^g}(X,Y)\simeq\mathrm{Hom}_{W_S}(X,Y)$, mais ça résulte de la formule d'adjonction et de $Y\to C(Y)$ dans $W_S^d\cap W_S^g$ pour $Y$ dans $\underline{\mathrm{Fibhot}}\,S$.}

où $\mathrm{Hom}_{W_S^d}$, $\mathrm{Hom}_{W_S^g}$ et $\mathrm{Hom}_{W_S}$ désignent respectivement les \deleted{\unsure{ens.}} $\mathrm{Hom}$ dans les catégories localisées par $W_S^d$, $W_S^g$ et $W_S$, i.e. dans $\mathrm{HOT}_W(S)$, $\mathrm{HOT}_W^{\unsure{\circ|o}}(S)$ et $\mathrm{HOT}_W^{\unsure{\mathrm{lc}|\mathrm{le}}}(S)$ respectivement.

\uline{Commentaire} Toutes les constructions et les énoncés de \textcircled{I} à \textcircled{VI}, à l'exception du fait que les foncteurs de descente induisent des équiv. pour les localisés \uline{grossiers} (pour \deleted{\unsure{les $S$}} les localisés grossiers), \inserted{et des formules d'adjonction de \textcircled{IV},} sont de nature quasiment soritale, et n'exigent d'autre \deleted{\unsure{propr.}} condition sur le localiseur \unsure{fondt|fondtl} $\unsure{W|\mathcal{W}}\subset\mathrm{Fl}\,\mathrm{Cat}$ que \deleted{\unsure{les celles de}} L1, L3 \unsure{faible}, L4 bis, récapitulées comme a) b) c) de la page 45. (On n'a même pas à utiliser d).) Il est vrai qu'en l'absence de la condition L5 (page 43), on ignore si à l'étage les deux localiseurs pertinents, $W_S$ et $W_S^f$ (égal à $W_S^d$ du côté gauche, à $W_S^g$ du côté droit) sont égaux. Donc à priori on trouve à l'étage \uline{deux} catégories localisées \unsure{évt|cert.|a priori} distinctes, et non pas une seule, savoir $(\underline{\mathrm{Fibhot}}\,S)\Sigma^{-1}$ avec $\Sigma=W_S$ ou $\Sigma=W_S^f$ $(=W_S^d=W_S^g)$.

Par contre, le fait que les foncteurs de descente induisent des équivalences entre les localisés grossiers, est plus profond, il exige \deleted{\unsure{la f}} \inserted{(me semble-t-il)} qu'on connaisse un foncteur ``montant'' $C$, disons de $\mathrm{Cat}_{/S}$ vers $\underline{\mathrm{Fibhot}}\,S$, \deleted{\unsure{ayant transformant les}} compatible avec $W_S$, et muni de

\[
X\xrightarrow{i_X}C(X)
\]

qui soit de $W_S$ pour $X$ quelconque. Quand on sait que $\mathrm{Cat}$ est une cat. de Quillen, même sans avoir les propriétés du foncteur $\underline{\mathrm{Ch}}_\infty$ ou $\underline{\mathrm{Ch}}$,\textsuperscript{*} on doit sans doute

% slanted
\margin{* \unsure{ni|ou} même être sûr de l'existence d'une factorisation \uline{fonctorielle} de $X\xrightarrow{f}S$ en $X\xrightarrow{i_X}C(X)\xrightarrow{\bar f}S$, \unsure{où} $i_X\in W$ et $\bar f$ \unsure{parf.}}
\scanpage{137}
\mspage{82}
% archive sheet number: 136

pouvoir s'en débrouiller. Mais la formule d'adjonction me semble rester inattendue et profonde, même en admettant les résultats de Quillen \inserted{et} Thomason. Elle paraît liée \inserted{ici} à \deleted{une} propriété assez particulière du foncteur $C$, savoir le lemme p.~68

% slanted
\margin{Voir \unsure{après coup|ap. ci-après} le th. sch. \unsure{2, -- 2°}, p 108.113. Ce n'est finalement pas si prof\unsure{ond} qu'il m'en semblait, mais quasiment sorital!}

\[
\begin{aligned}
,\,i_{C(X)}&=C(i_X):C(X)\unsure{\mathrel{\relbar\joinrel\relbar\joinrel\vert}|\longrightarrow}\unsure{C^2(C(X))|C(C(X))}\\
&\overset{\uparrow}{\text{dans }}\underbrace{(\mathrm{Cat}_{/S})(W_S^d)^{-1}}_{\mathrm{HOT}_W(S)}
 \text{ et dans }\underbrace{(\mathrm{Cat}_{/S})(W_S^g)^{-1}}_{\mathrm{HOT}_W^{\unsure{\circ|o}}(S)},
\end{aligned}
\]

ce qui n'a pas l'air purement formel. \deleted{\unsure{Il le semble}} Il serait intéressant d'examiner dans quelle mesure \deleted{elle} \deleted{\unsure{pourrait}} cette formule garde un sens \deleted{pour} dans une ``catégorie de modèles'' $(\mathcal M,W)$ plus ou moins quelconque, \deleted{\unsure{du moins en}} ou provenant d'un triple de Quillen $(W,\mathrm{Cof},\mathrm{Fib})$ dans $\mathcal M$. Il est clair que la factorisation $X\overset{f}{\unsure{\mathrel{\relbar\joinrel\relbar\joinrel\vert}|\longrightarrow}}S$ en $X\overset{i_X}{\unsure{-|\longrightarrow}}C(X)\overset{\bar f}{\unsure{-|\longrightarrow}}S$ ($i_X\in W$, $\bar f\in\mathrm{Parf}$, jouant le rôle des ``fibrations'' dans le contexte Quillen, a un sens intrinsèque indépendt du choix d'une construction particulière, chose qu'on peut préciser de façon parfaite. C'est la \deleted{\unsure{réd.}} description \unsure{des localiseurs|du localiseur} $W_S^d,W_S^g$ dans $\mathcal M_{/S}$ (pour $S\in\mathrm{Ob}\,\mathcal M$) qui n'est pas comprise de façon générale, en dehors du cas $\mathcal M=\mathrm{Cat}$. En d'autres termes, je ne comprends pas comment construire des catégories $\mathrm{\unsure{HOt|HOT}}_{\mathcal M,W}(S)$, $\mathrm{HOT}_{\mathcal M,W}^{\circ}(S)$, $=$ si $(\mathcal M,W)$ provient d'un triple de Quillen, j'\deleted{\unsure{ai|avais}} \unsure{en|n'}\inserted{vois} seulement la description d'une catégorie $\mathrm{HOT}^{\unsure{\mathrm{lc}|\mathrm{le}}}_{\mathcal M,W}(S)$ \inserted{$=(\mathcal M_{/S})(W_S^{\mathrm{grossier}})^{-1}$}, et un foncteur \inserted{de localisation}

\[
\mathcal M_{/S}\longrightarrow\mathrm{HOT}_{\mathcal M,W}^{\mathrm{lc}}(S).
\]

Il faut arriver à décrire des \inserted{$\mathcal M_{/S}$ \deleted{\unsure{des}}} $W$-équivalences (``\uline{relatives à $S$}'') plus fines que la grossière, \deleted{\unsure{induite}} \inserted{donnée} par $W$ sur la source $X$ sans plus. Dans $\mathrm{Cat}$ j'y arrive en utilisant \inserted{soit} la famille des isom. locaux (pour décrire $W_S^{\deleted{\unsure{g}}d}$), soit celle des isom colocaux (pour décrire $W_S^g$) -- en tous les cas il y a une notion de ``localisation'' (soit sur $S$, soit sur $S^\circ$\ldots), ou
\scanpage{138}
\mspage{83}
% archive sheet number: 137

\unsure{sort|dont} la richesse des intuitions liées à une telle notion. Peut-être, sans le savoir ai-je toujours traité $\mathrm{Cat}$ non comme une catégorie ``nue'', mais comme un site, ou plus \unsure{t\textsuperscript{t}|tt} (puisqu'il y a l'involution $X\to X^{\unsure{\circ|o}}$ qui ne peut être éludée), comme un ``bi-site'' -- on y a \uline{deux} manières de localiser, \deleted{\unsure{par}} soit par les $X_{/x}$ (top. de Zariski!) soit par les $\,{}_x\backslash X$ (top. co-Zariskienne). Il y a aussi la topologie lisse (toute proche sans doute de l'intuition de la \unsure{cat.} étale), et la topologie propre -- et on s'attend que la lisse donne ess. les $\unsure{\Rightarrow|=}$ \unsure{résultats} que la Zariskienne \deleted{\unsure{-- mais peut-être pas juste}} (vu que la topologie étale \deleted{\unsure{plus}} est égale -- il n'y a pas lieu d'en faire un plat\ldots) ; et la topologie propre les mêmes résultats (p.ex. la $=$ cohomologie ?) que la top. co-Zariskienne.

Il semble bien que dans \deleted{\unsure{toutes}} les catégories de modèles utilisées jusqu'à présent pour décrire $\mathrm{Hot}$ comme un localisé, on dispose de notions de localisation: $(\mathrm{Top})$, $\Delta^\wedge$ (ou $A^\wedge$, avec $A$ catégorie-test quelconque), $(\mathrm{Cat})$; et dans $\mathrm{Cat}$ on dispose même de \uline{deux} telles notions. Je pense par contre aux $\infty$-groupoïdes, c'est peut-être la seule catégorie de modèles à tel point proche déjà de $\mathrm{Hot}$ \unsure{elle-même}, qu'on n'y discerne pas de notion de localisation, et que l'idée même de \unsure{vouloir|modèles} \unsure{s'y localiser|s'y localise} (si ce n'est en ``localisant'' $\mathrm{Hot}$ lui-même) paraît quasiment
\scanpage{139}
\mspage{84}
% archive sheet number: 138

saugrenue.

Cette rapide digression me convainc à présent que la compatibilité dite pour la construction de $C(X)$ et $i_{\unsure{y|X}}:X\unsure{-|\to}C(\deleted{\unsure{A|M}}X)$, ne peut en aucun cas être purement formelle, une pure histoire de catégories de Quillen disons. \deleted{\unsure{Le}} \deleted{\unsure{je pourrais}} \inserted{\unsure{demander|du moins}} au mieux la compatibilité dans $\mathcal M_{/S}W_S^{-1}$, ce qui \deleted{est} ridicule! \deleted{On peut arguer} Mais si on prend comme localiseur $\Sigma$ dans $\mathcal M_{/S}$ les isomorphismes sans plus, la relation de compatibilité doit devenir fausse à tous les coups! Comme ensemble de flèches \inserted{$\Sigma$ un peu} plus raisonnable, qui soit \inserted{encore} très petit, on peut songer p.ex. à la $W$-équivalence \uline{universelle} sur $S$ (nettement plus fine encore que $W_S^d,W_S^g$ dans le cas particulier de $\mathrm{Cat}$). Aurait-on encore compatibilité pour cet ens. de flèches ? Je suis aussi un peu perplexe par le fait (pas bien compris) que les formules d'adjonction pour les topologie Zariskienne et coZariskienne \deleted{\unsure{soient de}} en sens opposé. (Ou me serai-je trompé en devinant \inserted{à la va-vite} ?) Mais à vrai dire, cette compatibilité à elle seule, même quand on utilise une structure de site sur $\mathcal M$, ne suffit pas à écrire une formule d'adjonction. On a deux catégories localisées,

\[
\mathrm{HOT}^{\unsure{\mathrm{lc}|\mathrm{le}}}_{\deleted{\unsure{M}}W}(S)=\mathcal M_{/S}W_S^{-1},\quad
\unsure{\mathcal M|\mathrm{d}}\mathrm{HOT}_{W}(S)=\mathcal M_{/S}(W_S^f)^{-1},
\]

où $W_S^f$ désigne une $W$-équivalence ``fine'', correspondant à la notion de localisation, plus exactement \unsure{aux|au} chgt. de base ``admissibles'' qu'on a choisis, et si \unsure{on} a un foncteur de localisation
\scanpage{140}
\mspage{85}
% archive sheet number: 139

\[
\mathrm{Hot}^{\deleted{\unsure{\mathrm{lc}}}}_W(S)\longrightarrow\mathrm{Hot}^{\unsure{\mathrm{lc}|\mathrm{le}}}_W(S),
\]

mais pas de foncteur en sens inverse. Mais si, mais si, \deleted{\unsure{je vois poindre l'aspect géom. de ma construction}, car les chgts de base (lisses disons) qui m'avaient servi pour définir \deleted{\unsure{l'équ}} la $W$-équivalence fine sur $S$, \unsure{me donnent un certain chgt}} en travaillant avec la catégorie similaire \deleted{\unsure{de}} à $\underline{\mathrm{Fibhot}}\,S$, ça doit \unsure{garder|prendre} un sens, tout au moins quand on a une bonne notion de ``fibration''. Je vois un petit travail de dégrossissage conceptuel à faire, avec peut-être une formule d'adjonction à la \unsure{``clef|clef}, mais je ne vais pas m'y lancer à présent.

$\backslash$

\margin{\unsure{V.16|Ve.16|VI.16}}

Je vais quand même essayer de faire ce travail ``de dégrossissage'', car le diagramme des 15 catégories a l'air un peu touffu, finalement! Je me donne

\[
(\mathcal M,W)\qquad W\subset\mathrm{Fl}\,\mathcal M,
\]

un $S\in\mathrm{Ob}\,\mathcal M$, d'où $\mathcal M_{/S}$, avec $W_S\subset\mathrm{Fl}\,(\mathcal M_{/S})$ induit par $W$. On va poser

\[
\mathrm{Ho}^{\mathrm{lc}}(S)=(\mathcal M_{/S})W_S^{-1},
\]

ça jouera le rôle \unsure{du type d'homotopie relativement constant|des types d'homotopie relativement constants}. Je suppose d'autre part que l'on a distingué une classe $\unsure{\mathcal F|F}\subset\mathrm{Fl}\,(\mathcal M)$ de flèches, jouant le rôle de ``fibrations'' (ce seront les morphismes Hot-fibrants de $\mathrm{Cat}$, donnant lieu à la s.catégorie $\underline{\mathrm{Fibhot}}\,S$ de $\mathrm{Cat}_{/S}$). Je prévois qu'il faudra
\scanpage{141}
\mspage{86}
% archive sheet number: 140

supposer, outre que $\mathcal M$ est stable par $\underset{\leftarrow}{\lim}$ finies, que $\unsure{F|\mathcal F}$ est stable par changement de base quelconque, et stable par composition, qu'il contient les iso. Et surtout qu'on a un th. de factorisation

\diagram{\begin{tikzcd}[ampersand replacement=\&,column sep=1.1em,row sep=0.5em]
X\arrow[rr,"f"]\arrow[dr,"i_X"']\&\& S\& i_X\in W,\;\bar f\in F\\
\&\bar X\arrow[ur,"\bar f"']\&\&
\end{tikzcd}}

pour toute flèche $f$ dans $\mathcal M$. Chez Quillen, on \deleted{exige} même que $i_X$ ait la propriété de relèvement à droite \deleted{pour les} relativement à $F$, mais ça ne devrait pas être essentiel, puisque je n'ai pas eu à utiliser ce fait dans $\mathrm{Cat}$. Quand \deleted{\unsure{je}} travaille dans la catégorie $\mathcal M_{/S}$, où $f$ est \unsure{sous-entendus|sous-entendue|sous-entendu} pour un objet $X$ de ladite, je noterai (par abus de langage) $C(X)$ pour $\bar X$, bien qu'il n'est pas dit qu'on ait un choix canonique de $C(X)$, encore moins que $C(X)$ dépende fonctoriellement de $X$.

Désignons par

\[
\underline{\mathrm{Fib}}\,S\subset\deleted{\unsure{\mathrm{C}}}\mathcal M_{/S}\tag*{(\ensuremath{\unsure{a|\alpha}})}
\]

la sous-catégorie de $\mathcal M_{/S}$ formée des $X$ dont le morphisme structural $f_X:X\unsure{-|\to}S$ est de $F$. \deleted{\unsure{On a}} Désignons \inserted{(par abus de notation)} par $W_S$ le localiseur induit sur $\underline{\mathrm{Fib}}\,S$ par $W_S$, donc l'inclusion \unsure{($*$)|(a)} définit

\[
(\underline{\mathrm{Fib}}\,S)W_S^{-1}\xrightarrow{\delta}(\mathcal M_{/S})W_S^{-1}\overset{\mathrm{def}}{=}\mathrm{Ho}^{\unsure{\mathrm{lc}|\mathrm{le}}}_W(S).
\]
\scanpage{142}
\mspage{87}
% archive sheet number: 141

Une première question, c'est de voir que c'est une équivalence. On voit que c'est ess. surjectif, \unsure{grâces|grâce} à l'existence de $C(X),i_X$. C'est peu! On va essayer de définir une flèche en sens inverse

\[
\bar C:\mathcal M_{/S}W_S^{-1}\unsure{\mathrel{\relbar\joinrel\relbar\joinrel\vert}|\longrightarrow}(\underline{\mathrm{Fib}}\,S)W_S^{-1},
\]

i.e. un \uline{foncteur}

\[
\gamma:\mathcal M_{/S}\longrightarrow(\underline{\mathrm{Fib}}\,S)W_S^{-1}
\]

qui transforme $W_S\subset\mathrm{Fl}\,(\mathcal M_{/S})$ en iso. On choisira, pour $X$ \unsure{ds|dans|de} $\mathcal M_{/S}$, un $C(X)$ dans $\underline{\mathrm{Fib}}\,S$ et un

\[
i_X:X\longrightarrow C(X)\qquad \unsure{v_X|i_X}\in W_S,
\]

et on posera

\[
\gamma(X)=C(X).
\]

La difficulté, c'est de montrer que ça ne dépend pas des choix faits, à isom can. près, et que c'est fonctoriel en $X$ -- et le deuxième implique le premier.

Pour voir que ça marche, visiblement il faut supposer quelque chose de plus. Je vois deux hypothèses qui, l'une et l'autre, \unsure{pourraient} faire marcher: l'une, c'est qu'on \deleted{\unsure{se}} donne déjà $C$ comme foncteur \inserted{et $i_X$ fonctoriel en $X$} (c'est ce qui était le cas dans $\mathrm{Cat}$ \unsure{tantôt|de fait}). L'autre, c'est que
\scanpage{143}
\mspage{88}
% archive sheet number: 142

les $i_X$ ont, \unsure{prr.|par.} à $\unsure{F|\mathcal F}$, la propriété de relèvement à la Quillen.

Je commence par l'hyp. que $C^{\unsure{2|*|?}}$ est déjà un foncteur. Le diagramme

\diagram{\begin{tikzcd}[ampersand replacement=\&]
X\arrow[r,dash,"i_X"]\arrow[d,"u"']\&CX\arrow[d,"C(u)"]\\
X'\arrow[r,dash,"i_{X'}"]\&CX',
\end{tikzcd}}

où $i_X,i_X\in W$, montre que $C$ transforme $W_S$ en $W_S$, donc en des iso de $\mathrm{Ho}^{\unsure{\mathrm{lc}|\mathrm{le}}}_W(S)$. Donc on tient le foncteur cherché $\unsure{\sigma|\gamma}$. On a en fait des foncteurs

\[
\deleted{\mathrm{id}\xrightarrow{\sim}\delta\gamma}
\]

\deleted{à l'aide des}

\diagram{\begin{tikzcd}[ampersand replacement=\&,column sep=large]
\underline{\mathrm{Fib}}\,S\arrow[r,hook,"{\overset{\mathrm{incl.}}{D}}"]\&\deleted{\unsure{\mathrm{C}}}\mathcal M_{/S}\arrow[l,bend left=18,"C"]
\end{tikzcd}}

compatibles avec les localiseurs $W_S$, induisant $\delta$ et $\sigma$, et on a

\[
\mathrm{id}_{\mathcal M_{/S}}\longrightarrow\deleted{\unsure{C}}DC,\qquad\text{de }W_S\ \text{\unsure{obj. par obj.|arg. par arg.}},
\]

\deleted{\unsure{car}} on a \deleted{\unsure{iso}}

\[
X\xrightarrow{i_X}DC(X)=C(X)
\]

avec $i_X\in W_S$. De $=$, on a

\[
\mathrm{id}_{\underline{\mathrm{Fib}}\,S}\longrightarrow CD\qquad\text{de }W_S\ \text{\unsure{obj. par obj.|arg. par arg.}},
\]

\unsure{avec} le même

\[
X\xrightarrow{i_X}CX\qquad(\text{cette fois }X\text{ de }\underline{\mathrm{Fib}}\,S),
\]
\scanpage{144}
\mspage{89}
% archive sheet number: 143

Donc on a sans larmes les deux foncteurs $\delta,\unsure{\gamma|\sigma}$ comme quasi-inverses l'une de l'autre. Et le lemme page 11 montre même que pour \uline{\unsure{tte|toute}} sous-catégorie pleine $\unsure{\mathcal F|F}$ de $\deleted{\unsure{\mathrm{C}}}\mathcal M_{/S}$, telle que

\[
\underline{\mathrm{Fib}}\,S\subset\mathcal F\subset\mathcal M_{/S},
\]

les inclusions induisent des équivalences sur les localisées $W_S$

\diagram{\begin{tikzcd}[ampersand replacement=\&,column sep=1.25em,row sep=0.4em]
(\underline{\mathrm{Fib}}\,S)W_S^{-1}\arrow[r,-{Bar},"\approx"]\&\mathcal F W_S^{-1}\arrow[r,-{Bar},"\approx"]\&\underline{(\mathcal M_{/S})W_S^{-1}}\arrow[d,equal,"\mathrm{def}"]\\
\&\&\mathrm{Ho}^{\unsure{\mathrm{l.c}|\mathrm{l.e}}}_W(S).
\end{tikzcd}}

Je vais traiter maintenant le cas du type Quillen, pour me convaincre que l'existence du foncteur $C$, et \unsure{les choix particuliers|le choix particulier} qu'on peut en faire, est \unsure{irrelevant|irrelev.nt}. (Ce dernier point est d'ailleurs évident à posteriori, quand on sait que $\delta$ est une équivalence, car le qu.inv. $\gamma$ sera alors unique à isom unique près, et le choix d'un foncteur $C$, et d'un \unsure{hom|lien} fonctoriel $i:\mathrm{id}\unsure{-|\to}C$, n'est qu'une façon parmi d'autres d'expliciter le quasi-inverse.)
\scanpage{145}
\mspage{90}
% archive sheet number: 144

\unsure{.} On se place donc dans l'hyp. quillénienne sur les $i_{\deleted{\unsure{Y}}X}$. On choisit $C(X),i_X$ pour chaque $X$ \inserted{dans $\mathcal M_{/S}$}, et on pose comme tantôt

\[
\gamma(X)=C(X)\qquad\text{comme objet de }(\mathcal M_{/S}\unsure{,|)}W_S^{-1}.
\]

Montrons que c'est fonctoriel en $X$. Soit

\[
u:X\unsure{-|\to}X'
\]

dans $\mathcal M_{/S}$, considérons le diagramme anodin \inserted{(\unsure{comm.})}

\diagram{\begin{tikzcd}[ampersand replacement=\&,column sep=1.4em,row sep=1.3em]
\deleted{\unsure{Y}}X\arrow[rr,"u"]\arrow[d,"i_X"']\arrow[drr,dashed,"{v=i_{X'}u}"]\&\&X'\arrow[d,"i_{X'}"]\\
C(X)\arrow[rr,dotted,"\bar u"]\arrow[dr]\&\& C(X')\arrow[dl,"{\in F}"]\\
\&S\&,
\end{tikzcd}}

\margin{Lemme: soit
\diagram{\begin{tikzcd}[ampersand replacement=\&,column sep=0.6em,row sep=0.55em,font=\scriptsize]
X\arrow[r,-{Bar},"u"]\arrow[d,dash,"i"']\&X'\arrow[d]\\
C\arrow[r,dotted]\&C'
\end{tikzcd}}
dans $\mathcal M_{/S}$ avec $i\in\mathrm{Cof}^{\deleted{\unsure{a}}}_0$ et $(C'\unsure{-|\to}S)\in\unsure{F|\mathcal F}\deleted{\unsure{=\mathrm{Fib}}}$ \inserted{$(=\mathrm{Fib})$}, alors la flèche en pointillés existe, \inserted{dans $\mathcal M_{/S}$}, de façon à rendre le carré comm.}

complétons le par la flèche en \inserted{gros} pointillés, $v=i_{X'}u$ (rendant le triangle sup. commutatif), l'hypothèse sur $i_X$ implique que la flèche $\bar u$ en pointillés fins existe, de façon à avoir encore deux triangles commutatifs. Mais ça signifie exactement qu'il y a un \uline{$S$-morphisme}, $\bar u$, i.e. une flèche dans $\deleted{\unsure{\mathrm C}}\mathcal M_{/S}$, rendant commutatif le carré

\diagram{\begin{tikzcd}[ampersand replacement=\&]
X\arrow[r,"u"]\arrow[d,dash,"i_X"']\& X'\arrow[d,"i_{X'}"]\\
C(X)\arrow[r,"\bar u"]\&C(X'),
\end{tikzcd}}
\scanpage{146}
\mspage{91}
% archive sheet number: 145

et il s'impose de prendre

\[
\gamma(u)=\text{image de }\bar u\text{ dans }\operatorname{Fl}\bigl(\underline{\mathrm{Fib}}\unsure{(\underline S)|(\underline S)_{S}}W_S^{-1}\bigr),
\]

Il faut prouver que cette flèche ne dépend pas du choix de $\bar u$. On doit dégager un lemme, où on désigne par \deleted{$\unsure{\mathrm{F}|\mathrm{Fi}}$} $\mathrm{Cof}_0$ l'ens des flèches de $\unsure{\mathcal M|M}$ qui ont la prop. de \unsure{relèvement|relèvemt} par. aux éléments de $\mathrm{Fib}=F$. ($\mathrm{Cof}_0$ joue le rôle des cofibrations triviales.)

\uline{Lemme} Soit un diagramme en traits pleins dans $\mathcal M_{/S}$

\diagram{\begin{tikzcd}[ampersand replacement=\&]
X \arrow[r,"u"] \arrow[d,"i"'] \& X' \arrow[d,"i'"]\\
C \arrow[r,dotted,"v"] \& C',
\end{tikzcd}}

avec $i\in\mathrm{Cof}_0$, et $C'\in\operatorname{Ob}\underline{\mathrm{Fib}}\,S$ i.e. $(\deleted{\unsure{S|B}}C'\unsure{\relbar|\longrightarrow} S)\in\mathrm{Fib}$. Alors $v$ existe, de façon à rendre le carré commutatif. De plus, \inserted{\deleted{si $C\in\operatorname{Ob}\underline{\mathrm{Fib}}\,S$ \unsure{et|si} \illegible{3}} si $i\in W$ et $C\in\operatorname{Ob}\underline{\mathrm{Fib}}\,S$,} le choix de $v$ \deleted{\unsure{dans}} comme \deleted{\unsure{l'ens. des}} flèche de $(\underline{\mathrm{Fib}}\, S) W_S^{-1}$ ne \deleted{dépend pas} \inserted{est indépendant} du choix fait.

Il reste à prouver cette indépendance.

\deleted{Supposons que \unsure{nous|l'on} ayons deux tels choix $(v_1,v_2)$, d'où}

\[
\deleted{v=(v_1,v_2):C\longrightarrow C'\times_S C',\ \deleted{\unsure{r|s}}}
\]

\deleted{rendant commutatif le diagramme}

\diagram{\deleted{\begin{tikzcd}[ampersand replacement=\&]
X \arrow[r,"u"] \arrow[d,"i"'] \& X' \arrow[d,"{(i',i')=\unsure{j|J}}"]\\
C \arrow[r,"v"] \& C'\times_S C'
\end{tikzcd}}}
